\documentclass[a4paper,12pt]{article}
\usepackage[utf8]{inputenc}
\usepackage[T1]{fontenc}
\usepackage[margin=2.5cm]{geometry}
\usepackage{ctex} % 支持中文
\usepackage{setspace} % 设置行距
\usepackage{titlesec} % 标题格式
\usepackage{xcolor} % 颜色支持

% 设置行距为1.5倍，增加阅读舒适度
\onehalfspacing

% 设置标题格式
\titleformat{\section}{\Large\bfseries\color{blue!40!black}}{}{0em}{}
\titleformat{\subsection}{\large\bfseries\color{gray!80!black}}{}{0em}{}

% 元数据
\title{\textbf{2024年度环境、社会及管治（ESG）报告节选}}
\author{公司董事会办公室}
\date{2024年3月}

\begin{document}

\maketitle

\section{1. 前言：服务国家战略，守护自贸港门户}

2024年，对于我们而言，是极不平凡的一年。这一年，是全省推进自贸港封关运作的攻坚之年，也是我们作为本土航空基础设施平台迎接新机遇的关键之年。本公司始终牢记“外防输入、内保畅通”的使命，紧紧围绕“一主一次，两支一货运”的机场群战略布局，全力提升枢纽能级。

我们深知，机场不仅是飞机的起降地，更是连接世界的桥梁。在服务国家开放战略的同时，我们也在思考如何让企业发展更可持续。这一年，我们不仅关注旅客吞吐量和航班起降架次，更将目光投向了那些看不见但至关重要的领域——如何减少每一吨碳排放，如何保障每一次起落的绝对安全，以及如何让每一位旅客感受到“家”的温暖。

\section{2. 环境范畴：守护碧海蓝天}

绿色，是自贸港最鲜明的底色。我们坚持“生态优先、绿色发展”的理念，通过一系列技术改造和精细化管理，努力降低运营过程中的环境足迹。

\subsection{2.1 我们的“碳账本”与能源管理}

为了让公众更清晰地了解我们的环境表现，我们对2024全年的排放数据进行了详细盘点。在温室气体排放方面，本年度公司排放总量控制在 51,518.71 吨二氧化碳当量。如果细分来看，直接排放（范围1，主要来自车辆汽油燃烧等）为 11,214.39 吨；间接排放（范围2，主要来自外购电力）为 40,304.32 吨。折算下来，我们的碳排放强度为 0.12 吨二氧化碳当量/万元。

在能源资源消耗方面，我们坚持“精打细算”。全年综合能源消耗量折合标准煤 13,538.22 吨。其中，外购电力消耗为 7,511.05 万千瓦时，汽油消耗为 19.10 万升。水资源管理同样是我们关注的重点，全年水资源消耗总量为 145.94 万吨。值得一提的是，我们非常重视水资源的循环利用，全年实现了 13.61 万吨的中水回用，让珍贵的水资源得到了二次生命。此外，在废弃物处理上，我们产生的 6,230.58 吨无害废弃物和 0.58 吨有害废弃物，全部实现了 100\% 的合规处理。

\subsection{2.2 打造“双碳机场”：让飞机在地面“静音”休息}

在机场的日常运行中，有一个往往被忽视的污染源，那就是飞机停在地面时使用的辅助动力装置（APU）。通俗地说，以前飞机停稳后，为了维持舱内的空调和照明，需要一直开着自带的小发动机烧油发电，这不仅噪音大，而且排放高。

为了解决这个问题，我们在旗下核心门户机场大力推广了“APU替代设施”。这项技术就像是给飞机配了一个巨大的“充电宝”。当飞机停靠廊桥时，直接连接地面的供电和供气设备，飞机就不再需要消耗航空燃油。

这一改变带来了显著的环保效益。2024年，我们共投入使用了31套APU替代设备，其中近机位18套，远机位13套。通过这一举措，我们全年累计减少了碳排放 25,382 吨二氧化碳当量。凭借在这一领域的出色表现，旗下核心门户机场荣获了行业内极具分量的“双碳机场”二星级称号。

\subsection{2.3 这里的阳光能发电：光伏与绿色建筑}

海南拥有得天独厚的阳光资源，我们充分利用这一优势，让屋顶变成了绿色的“发电厂”。在旗下核心商业综合体，我们利用闲置屋顶建设了光伏发电项目。这个项目年均发电量达到 450 万千瓦时，这些清洁电力直接用于商业体的日常运营，占到了其总用电量的 10\%，相当于每年减少碳排放 2,263.59 吨。在旗下的临空产业园，我们更是铺设了装机容量达 8.6MWp 的屋顶光伏，预计每年可减少碳排放 25.75 万吨。

除了生产绿电，我们还在用车和建筑上做“减法”。在核心机场，我们积极推动车辆“油改电”，目前新能源车占比已提升至 43.6\%，总量达到 221 辆，并配套新建了约 100 个充电桩，全年因此减碳 761 吨。在建筑领域，正在建设中的省会城市地标项目“海南中心”（双子塔）已经获得了 LEED CS 金级预认证，而某新区的 F07 项目也达到了绿色建筑一星级标准。

\section{3. 社会范畴：共建美好空港}

机场是城市的窗口，也是社会的缩影。2024年，我们用实际行动诠释了什么是“真情服务”，什么是“责任担当”。

\subsection{3.1 运营复苏与服务细节的温度}

随着旅游市场的全面复苏，我们的业务数据也划出了一条漂亮的上扬曲线。全年旗下控股及管理输出的9家机场总体起降架次达到 16.11 万架次，旅客吞吐量共计 2,462.62 万人次。特别是我们的旗下核心门户机场，旅客吞吐量成功重回“2000万级”俱乐部，位列全国第29位。与此同时，旗下重点区域机场也迎来了历史性突破，开通了至吉隆坡的首条国际客运航线。

但在这些宏大的数据背后，我们更看重每一位旅客的体验。为了让通关更顺畅，我们实施了“一机双屏”海关安检融合查验模式。过去，旅客需要分别排队过海关和安检，现在只需一次过检，行李图像同时呈现在海关和安检人员面前，大大缩短了等待时间。

针对外籍旅客入境后可能遇到的“支付难、办卡难”问题，我们在核心机场启用了“一站式综合服务中心”。这是一个充满温情的小角落，外籍友人在下飞机的第一时间，就能在这里解决移动支付绑定、购买电话卡、兑换货币等一系列生活琐事。我们希望用这种“落地即融入”的服务，向世界展示自贸港的热情与包容。

\subsection{3.2 直面风雨：台风“摩羯”下的坚守}

2024年9月，超强台风“摩羯”正面袭击本省，给我们的运营带来了前所未有的挑战。这不仅是一场自然灾害，更是一场对我们应急体系和团队意志的极限测试。

台风登陆前，我们迅速启动了“防汛防风专项应急预案”。运管委通宵达旦地进行航班动态调整，确保每一架飞机都能安全转场或系留。在机坪上，工作人员对所有设施进行了地毯式的加固，严阵以待。

当台风过境后，我们展现出了令人惊叹的“复航速度”。全体员工不分昼夜地清理跑道异物，修复受损的围界和导航设施。仅仅在9月3日至7日这短短几天内，核心机场就迅速恢复了常态化运行，累计执行航班 1,300 架次，为全省的救灾物资运输和人员进出岛搭建了最可靠的空中通道。

在日常安全管理上，我们也从未松懈。2024年，我们上线了“安全智控”平台，将厚厚的纸质手册变成了数字化的指令。平台关联了 9,626 项制度手册，全年生成了 31 万份电子检查单，确保每一个安全动作都有据可查。得益于此，全年我们实现了一般征候万架次率为 0，责任事故为 0，安全培训覆盖率 100\% 的优异成绩。

\subsection{3.3 员工成长与社会回馈}

员工是公司发展的基石。截至2024年底，大家庭共有员工 10,611 人，其中女性员工 4,133 人，占比 38.9\%；少数民族员工 179 人。我们非常重视员工的成长，全年培训总投入达到 1,119.54 万元，人均培训时长高达 105.88 小时。在公司治理层面，董事会中女性成员占比达到 22\%，体现了我们在性别平等方面的努力。

在谋求自身发展的同时，我们没有忘记回馈社会。2024年，我们投入乡村振兴资金 190 万元，帮助对口帮扶村改善基础设施和教育条件。我们的“红马甲”志愿者团队也活跃在各个角落，全年累计参与志愿服务超过 10,000 人次，总时长达到 500,000 小时。此外，我们还圆满完成了博鳌亚洲论坛年会、消博会、少数民族运动会等一系列国家级重大活动的保障任务，用专业的服务擦亮了“中国服务”的金字招牌。

\section{4. 管治范畴：筑牢合规基石}

透明、高效的治理是企业行稳致远的保障。我们严格按照上市公司标准，构建了由股东大会、董事会、监事会及管理层组成的“五会一层”治理体系。

本届董事会共有9名成员，其中独立董事3名，这种结构有效地保证了决策的客观性与独立性。2024年，我们共召开了5次股东大会和13次董事会，发布了70份各类公告，确保每一位投资者都能及时、公平地获取公司信息。

在风险控制方面，我们构建了“三全三化”大监督体系，全年梳理出358个潜在风险点并逐一制定防控措施。我们深知廉洁是底线，因此对管理层和员工分别实现了 95\% 和 90\% 的反腐败培训覆盖率，培训总时长达 400 小时。对于合作伙伴，我们同样严格要求，供应商廉洁承诺书签署率达到了 100\%，特别是在工程类供应商遴选中，严格执行合规标准。

最后，在经济绩效方面，作为自贸港建设的积极参与者，我们通过参股等形式布局离岛免税业务。2024年，公司间接参与的5家离岛免税店线下销售额约为 52 亿元，市场占比约 17\%，为促进消费回流、繁荣地方经济做出了实实在在的贡献。

\end{document}